\documentclass[11pt,a4paper]{article}
\usepackage[T1]{fontenc}
\usepackage[utf8]{inputenc}
\usepackage[a4paper,margin=1in]{geometry}
\usepackage{mathptmx}
\usepackage{setspace}
\usepackage{enumitem}
\usepackage[hidelinks]{hyperref}
\usepackage{titlesec}
\usepackage{parskip}

\setstretch{1.15}
\setlist[itemize]{leftmargin=*,itemsep=0.2em,topsep=0.3em}
\titleformat{\section}{\normalfont\large\bfseries}{\thesection}{1em}{}

\begin{document}

\begin{center}
{\large\bfseries Related Paper Review 1}\\[0.4em]
{\large Human-In-the-Loop Software Development Agents (HULA)}
\end{center}

\vspace{0.6em}

\begin{center}
\begin{tabular}{@{}ll@{}}
\textbf{Project} & A.R.T.E.M.I.S (F26-008), FAST NUCES Lahore \\
\textbf{Team} & Hashir Rauf (23L-2572), Hira Khalid (23L-2594), \\
 & Syeda Doureesha Batool (23L-2651) \\
\textbf{Supervisor} & Mr.\ Mubasher Hussain \\
\end{tabular}
\end{center}

\vspace{0.4em}
\hrule
\vspace{0.8em}

\section{Citation and Publication Status}

W. Takerngsaksiri, J. Pasuksmit, P. Thongtanunam, C. Tantithamthavorn, R. Zhang,
F. Jiang, J. Li, E. Cook, K. Chen, and M. Wu, ``Human-In-the-Loop Software
Development Agents,'' in \textit{Proceedings of the 2025 IEEE/ACM 47th
International Conference on Software Engineering: Software Engineering in
Practice (ICSE-SEIP 2025)}, Ottawa, Canada, pp. 342--352, IEEE, 2025.
DOI: 10.1109/ICSE-SEIP66354.2025.00036

This is a peer-reviewed paper published in the proceedings of ICSE, the leading
conference in software engineering. It is not a preprint. The work was carried
out and deployed at Atlassian, so the results come from an industrial
deployment rather than a laboratory study.

\section{What the Paper Does}

The paper presents HULA, a framework in which a large language model agent
writes software under human supervision rather than autonomously. The authors
observe that fully autonomous coding agents produce changes that engineers will
not merge, because the engineer has no opportunity to correct the agent's
intent before code is written.

HULA therefore splits the work into three stages with a person in control at
each boundary. An AI planner agent reads the work item and the repository and
produces a coding plan, naming the files it intends to change. The engineer
reviews that plan and may edit or reject it before anything is generated. An AI
coding agent then implements only the approved plan, and the engineer reviews
the resulting code before it becomes a pull request.

The framework was deployed inside Atlassian and integrated with Jira, so that
work items could be handed to the agent directly. The evaluation covers both
offline benchmarks and real usage by Atlassian engineers.

\section{Results Reported}

\begin{itemize}
  \item Engineers approved the agent's coding plan in approximately 82 percent
        of cases, indicating that a plan is a reviewable artefact that people
        can meaningfully judge before code exists.
  \item A substantially smaller proportion of generated pull requests were
        merged than plans approved, which shows that approval at the plan stage
        does not by itself guarantee acceptable code.
  \item Engineers reported that the framework reduced the effort of starting a
        task, particularly in forming an initial plan and writing code for
        straightforward work items.
  \item Output quality depended heavily on the quality of the context supplied
        to the agent, with insufficient context being a recurring cause of
        unusable output.
\end{itemize}

\section{Strengths}

The central strength is that the human-approval loop is validated in a real
organisation rather than argued for in principle. The 82 percent plan approval
rate is evidence that asking a person to approve a plan is a workable
interaction, not an obstacle that users route around.

The separation of planning from generation is the second strength. By making
the plan an explicit artefact that exists before any code is written, the
framework gives the user a cheap point of intervention. Correcting a plan costs
far less than reviewing and repairing generated code.

The deployment context is itself a contribution. Because the system ran inside
a working engineering organisation and against real Jira work items, the
reported friction is real friction rather than anticipated friction.

\section{Weaknesses and Limitations}

The scope is confined to software engineering tasks within a repository, and
within an organisation whose engineers are expert reviewers. The paper does not
address a non-expert user, who may be unable to judge whether a plan is
correct.

The gap between plan approval and pull request merge is acknowledged but not
resolved. A user may approve a reasonable plan and still receive code that
cannot be used, which means approval at one stage does not transfer to
confidence at the next.

The framework also has no reversal mechanism. Once a change has been applied,
the paper offers no facility to undo it as a unit, relying instead on the
surrounding version control system.

Finally, the agent operates with access to the repository as a whole. There is
no notion of a user-granted boundary restricting which files the agent may read
or modify.

\section{Relationship to A.R.T.E.M.I.S}

This paper is the closest published work to the interaction model used in our
project, and it supports a design decision we had already made. Our system also
produces an inspectable plan before acting, and holds that plan for the user to
approve, edit or reject. The 82 percent approval figure is external evidence
that this model is practical, which strengthens the case for our approval gate
beyond our own questionnaire finding.

Our project differs in four respects. First, our system applies the loop to
general desktop knowledge work rather than only to code, so the user is not
assumed to be a professional engineer. Second, our code assistance never writes
to a file at all: it emits a difference for the user to apply, whereas HULA
generates code into a branch. Third, we add reversal, so an approved plan can be
undone as the single unit that was approved. Fourth, and most importantly, our
system operates only inside folders the user has explicitly granted, so the
blast radius of an incorrect plan is bounded by construction rather than by the
agent's good behaviour.

The paper also supplies a warning we have adopted. Because output quality
tracked input context so strongly, our design surfaces the files and context
being used before proposing a change, rather than leaving the user to guess what
the system considered.

\end{document}
